\documentclass[../Mathe_Aufgaben.tex]{subfiles}
\begin{document}
\chapter{Grundlagen der Stochastik}

\section{Kombinatorik}
	\label{sec:kombinatorik}

\begin{komplexaufgaben}{sec:kombinatorik}

	\komplexaufgabe{
		Bearbeiten Sie die folgenden Aufgaben in mehreren Phasen:
		\begin{enumerate}[label=\Alph*)]
			\item Lösen Sie zunächst alle Teilaufgaben a) bis j).
			\item Sortieren Sie diese Aufgaben nach Aufgabentypen.
			\item Geben Sie Gesichtspunkte an, nach denen Sie Ihre Einteilung vorgenommen haben.
			\item Versuchen Sie die jeweilige kombinatorische Problemstellung und Ihren Lösungsweg allgemein zu formulieren.
			\item Lesen Sie nun Abschnitt 2.1 (Kombinatorik: Systematisches Abzählen) im Lehrtext und vergleichen Sie:
			\begin{itemize}
				\item Stimmen Ihre Kategorien mit den vier Grundfällen überein?
				\item Wie passen Ihre Aufgaben in die 2×2-Tabelle (mit/ohne Zurücklegen, mit/ohne Reihenfolge)?
				\item Welche Formeln können Sie auf Ihre Aufgaben anwenden?
			\end{itemize}
			\item Überprüfen Sie Ihre Lösungen aus Teilaufgabe A mit den Formeln aus dem Lehrtext.
		\end{enumerate}
	}{
		\item Wieviele sechsstellige Zahlen -- gebildet aus den Ziffern 1,...,9 -- gibt es?

		\item Von den 20 Schülern einer Klasse sollen 10 in die Parallelklasse überwechseln. Wieviele Möglichkeiten gibt es?

		\item Acht Athleten kämpfen um drei Medaillen. Wie viele Möglichkeiten für die Preisverteilung gibt es?

		\item Auf wieviele Arten kann man 5 Personen auf 8 Einzelzimmer verteilen?

		\item Herr M. will seine 5 Kinder für ein Gruppenphoto in einer Reihe anordnen. Wieviele Möglichkeiten hat er?

		\item Wieviele verschiedene Tippreihen gibt es bei der Elferwette im Fußballtoto (11 Tipps: jeweils gewonnen, verloren oder unentschieden)?

		\item In einem Raum gibt es 8 Lampen, die unabhängig ein- und ausgeschaltet werden können. Wieviele Möglichkeiten gibt es, so dass genau 5 Lampen brennen?

		\item Einem Skatspiel werden drei Karten mit Zurücklegen entnommen. Wieviele Möglichkeiten gibt es?

		\item Wieviele Möglichkeiten gibt es, 8 gleiche Gegenstände auf 12 Fächer zu verteilen?

		\item In einem Hotel sind noch 5 Zimmer frei. Jedes der Zimmer ist ein Dreibett-Zimmer. Am Abend kommen noch drei Gäste. Wie viele Möglichkeiten hat der Hotelier, die Gäste auf die Zimmer zu verteilen? Es wird dabei zugelassen, dass in einem Zimmer evt. 2 oder auch 3 Gäste schlafen.
	}

	\komplexaufgabe{Trapper Fuzzi muss auf dem Weg nach Alaska drei Flüsse überqueren.
		Am ersten Fluss gibt es 7 Furten, am zweiten 5 und am dritten 3 Furten.
		Am ersten Fluss sind 6, am zweiten 4 und am dritten 2 der Furten passierbar.
		Fuzzi entscheidet sich vor dem Abmarsch für eine Route.
		Sollte man darauf wetten, dass er durchkommt?}{}

	\komplexaufgabe{Eine Dualzahl hat nur Nullen und Einsen als Ziffern.
		Computer arbeiten mit 8-stelligen, 16-stelligen oder 32-stelligen Dualzahlen.
		Wie viele verschiedene Zahlen kann man für jede dieser Stellenanzahlen bilden?}{}

	\komplexaufgabe{Wie viele 7-stellige Zahlen ohne Ziffernwiederholung lassen sich aus den Ziffern 1 bis 9 bilden?}{}

	\komplexaufgabe{Bei einer Silvester-Party stößt um Mitternacht jeder der 7 Party-Gäste
		mit dem Sektglas mit jedem Gast an.
		Wie oft klingen dabei die Gläser?}{}

\end{komplexaufgaben}



\section{Klassische Wahrscheinlichkeit und Baumdiagramme}
\label{sec:klassWahrsch}

\begin{komplexaufgaben}{sec:klassWahrsch}

	\komplexaufgabe{Lösen Sie die folgenden Wahrscheinlichkeitsaufgaben.}{
		\item Aus einer Urne mit 15 weißen und 5 roten Kugeln werden 8 Kugeln ohne Zurücklegen gezogen.
			Mit welcher Wahrscheinlichkeit sind unter den gezogenen Kugeln genau 3 rote Kugeln?
			Mit welcher Wahrscheinlichkeit sind mindestens 4 rote dabei?
		\item In einer Lieferung von 100 Transistoren sind 10 defekt.
			Mit welcher Wahrscheinlichkeit werden bei Entnahme einer Stichprobe von 5 Transistoren
			genau 2 (mindestens 3) defekte Transistoren entdeckt?
	}

	\komplexaufgabe{Auf einem Rummelplatz wird ein Minilotto \glqq{}4 aus 16\grqq{} angeboten.
		Der Spieleinsatz beträgt pro Tipp $1\,\text{€}$.
		Die Auszahlungsquoten lauten $10\,\text{€}$ bei 3 Richtigen und $1000\,\text{€}$ bei 4 Richtigen.
		Mit welchem mittleren Gewinn kann der Veranstalter pro Tipp rechnen?}{}

	\komplexaufgabe{In einer Urne befinden sich 5 rote, 3 weiße und 6 schwarze Kugeln.
		3 Kugeln werden ohne Zurücklegen gezogen.
		Mit welcher Wahrscheinlichkeit sind sie alle verschiedenfarbig
		(alle rot, alle gleichfarbig)?}{}

	\komplexaufgabe{In einer Sendung von 80 Batterien befinden sich 10 defekte.
		Mit welcher Wahrscheinlichkeit enthält eine Stichprobe von 5 Batterien
		genau eine (genau 3, höchstens 4, mindestens eine) defekte Batterie?}{}

	\komplexaufgabe{Unter den 100 Losen einer Lotterie befinden sich 2 Hauptgewinne,
		8 einfache Gewinne und 20 Trostpreise.}{
		\item Mit welcher Wahrscheinlichkeit befinden sich unter 5 gezogenen Losen
			genau ein Hauptgewinn und sonst nur Nieten (überhaupt kein Gewinn)?
		\item Mit welcher Wahrscheinlichkeit befinden sich unter 10 gezogenen Losen genau
			2 einfache Gewinne, 3 Trostpreise und sonst nur Nieten
			(1 Hauptgewinn, 2 einfache Gewinne und sonst nur Nieten)?
			Anleitung: Teilen Sie die Lose in 3 geeignet gewählte Gruppen.
	}

	\komplexaufgabe{
		\textbf{Das Problem des Chevalier de Méré}

		Der französische Adlige Chevalier de Méré stellte 1654 folgende beiden Wetten vor:

		\textbf{Wette A:} „Ich werfe einen Würfel 4-mal. Ich wette, dass mindestens eine 6 fällt."

		\textbf{Wette B:} „Ich werfe zwei Würfel 24-mal. Ich wette, dass mindestens eine Doppel-6 (beide Würfel zeigen 6) fällt."

		De Méré glaubte intuitiv, beide Wetten seien gleich gut, verlor aber auf Dauer bei Wette B.

		\textbf{Arbeitsaufträge:}
	}{
		\item \textbf{Intuition:} Welche Wette würden Sie eingehen? Begründen Sie Ihre Intuition, \emph{ohne zu rechnen}.

		\item \textbf{Modellierung:} Beschreiben Sie für beide Wetten die Ergebnismenge $\Omega$, das gesuchte Ereignis sowie die Annahmen über die Würfe.

		\item \textbf{Berechnung Wette A:} Berechnen Sie $P(\text{keine 6 in 4 Würfen})$ und daraus $P(\text{mindestens eine 6})$.

		\item \textbf{Berechnung Wette B:} Bestimmen Sie $P(\text{Doppel-6 bei einem Wurf})$ und berechnen Sie $P(\text{mindestens eine Doppel-6 in 24 Würfen})$.

		\item \textbf{Vergleich:} Welche Wette ist günstiger? War Ihre Intuition aus a) richtig?

		\item \textbf{De Mérés Irrtum:} De Méré argumentierte: „$\frac{4}{6} = \frac{2}{3}$ bzw.\ $\frac{24}{36} = \frac{2}{3}$ Chance." Wo liegt der Fehler?

		\item \textbf{Historische Bedeutung:} Warum gilt dieser Briefwechsel zwischen Pascal und Fermat (1654) als Geburtsstunde der Wahrscheinlichkeitsrechnung?
	}

	% --- Grundlegende Bernoulli-Aufgaben ---

	\komplexaufgabe{Ein fairer Würfel wird viermal geworfen (Bernoulli-Kette mit $n=4$, $p=\tfrac{1}{6}$).}{
		\item Zeichnen Sie einen Ausschnitt des Baumdiagramms (erste zwei Stufen, Äste beschriften mit Wahrscheinlichkeiten).
		\item Berechnen Sie die Wahrscheinlichkeit, genau einmal eine 6 zu würfeln.
		\item Berechnen Sie die Wahrscheinlichkeit, mindestens einmal eine 6 zu würfeln.
			(Hinweis: Gegenereignis)
		\item Wie groß ist die Wahrscheinlichkeit, keine einzige 6 zu würfeln?
	}

	\komplexaufgabe{Eine faire Münze wird fünfmal geworfen.}{
		\item Berechnen Sie $P(\text{genau 3-mal Kopf})$.
		\item Berechnen Sie $P(\text{höchstens 2-mal Kopf})$.
		\item Berechnen Sie $P(\text{mindestens 4-mal Kopf})$.
	}

	\komplexaufgabe{Ein Basketballspieler trifft einen Freiwurf mit Wahrscheinlichkeit $p = 0{,}7$. Er wirft dreimal.}{
		\item Zeichnen Sie das vollständige Baumdiagramm (3 Stufen, Äste mit T für Treffer und F für Fehlwurf).
		\item Berechnen Sie die Wahrscheinlichkeit, genau zweimal zu treffen.
		\item Berechnen Sie die Wahrscheinlichkeit, mindestens einmal zu treffen.
	}

	\komplexaufgabe{Bei einem Fertigungsverfahren sind $5\,\%$ der Glühbirnen defekt.
		Es werden sechs Glühbirnen zufällig entnommen und geprüft.}{
		\item Berechnen Sie die Wahrscheinlichkeit, dass keine defekte Birne dabei ist.
		\item Berechnen Sie die Wahrscheinlichkeit, dass genau eine defekte Birne dabei ist.
		\item Berechnen Sie die Wahrscheinlichkeit, dass mindestens eine defekte Birne dabei ist.
	}

	\komplexaufgabe{Ein Multiple-Choice-Test besteht aus vier Fragen.
		Jede Frage hat drei Antwortmöglichkeiten, von denen genau eine richtig ist.
		Ein Schüler rät bei allen Fragen zufällig.}{
		\item Berechnen Sie die Wahrscheinlichkeit, alle vier Fragen richtig zu beantworten.
		\item Berechnen Sie die Wahrscheinlichkeit, genau zwei Fragen richtig zu beantworten.
		\item Berechnen Sie die Wahrscheinlichkeit, mindestens drei Fragen richtig zu beantworten.
	}

	\komplexaufgabe{Eine Schützin trifft die Scheibe mit Wahrscheinlichkeit $p = 0{,}8$. Sie schießt fünfmal.}{
		\item Berechnen Sie $P(\text{genau 4 Treffer})$.
		\item Berechnen Sie $P(\text{mindestens 4 Treffer})$.
		\item Berechnen Sie $P(\text{alle 5 Treffer})$ und vergleichen Sie die drei Ergebnisse.
	}

	% --- Parametervariation: unbekanntes n ---

	\komplexaufgabe{Ein fairer Würfel wird $n$-mal geworfen.}{
		\item Drücken Sie $P(\text{mindestens eine 6})$ als Funktion von $n$ aus.
		\item Wie oft muss man mindestens würfeln, damit $P(\text{mindestens eine 6}) \ge 0{,}90$?
			Lösen Sie die Ungleichung mit dem natürlichen Logarithmus.
	}

	\komplexaufgabe{Eine faire Münze wird $n$-mal geworfen.}{
		\item Berechnen Sie $P(\text{mindestens 3-mal Kopf})$ für $n = 4$ und für $n = 5$.
		\item Ab welchem $n$ gilt $P(\text{mindestens 3-mal Kopf}) \ge 0{,}50$? Überprüfen Sie durch Einsetzen.
	}

	\komplexaufgabe{Eine Maschine produziert Schrauben, von denen $10\,\%$ defekt sind.
		Es werden $n$ Schrauben geprüft.}{
		\item Geben Sie $P(\text{mindestens eine defekte Schraube})$ als Funktion von $n$ an.
		\item Wie viele Schrauben muss man mindestens prüfen, damit diese Wahrscheinlichkeit
			mindestens $90\,\%$ beträgt? Lösen Sie die Ungleichung mit dem Logarithmus.
	}

	% --- Parametervariation: unbekanntes p ---

	\komplexaufgabe{Bei vier unabhängigen Versuchen gilt $P(\text{genau 2 Treffer}) = \dfrac{216}{625}$.}{
		\item Stellen Sie eine Gleichung für die Einzeltrefferwahrscheinlichkeit $p$ auf und lösen Sie sie.
		\item Erklären Sie, warum es zwei Lösungen gibt.
	}

	\komplexaufgabe{Bei zwei unabhängigen Versuchen gilt
		$P(\text{mindestens 1 Treffer}) = \dfrac{5}{9}$.
		Bestimmen Sie die Einzeltrefferwahrscheinlichkeit $p$.}{}

	\komplexaufgabe{Bei drei unabhängigen Versuchen gilt
		$P(\text{kein Treffer}) = \dfrac{8}{27}$.
		Bestimmen Sie die Einzeltrefferwahrscheinlichkeit $p$.}{}

	% --- Parametervariation: unbekanntes k ---

	\komplexaufgabe{Ein fairer Würfel wird sechsmal geworfen ($p = \tfrac{1}{6}$).
		Berechnen Sie $P(\text{genau $k$ Sechsen})$ für $k = 0, 1, 2, 3$
		und bestimmen Sie, für welches $k$ diese Wahrscheinlichkeit am größten ist.}{}

	\komplexaufgabe{Eine faire Münze wird zehnmal geworfen.
		Für welche Werte von $k$ gilt $P(\text{genau $k$-mal Kopf}) \ge 0{,}20$?
		Berechnen Sie die relevanten Wahrscheinlichkeiten.}{}

	\komplexaufgabe{Bei einem Versuch mit Erfolgswahrscheinlichkeit $p = 0{,}3$ werden
		fünf unabhängige Versuche durchgeführt.
		Bestimmen Sie das kleinste $k$, für das
		$P(\text{mindestens $k$ Treffer}) \le 0{,}10$ gilt.}{}

\end{komplexaufgaben}


\section{Mengenoperationen und Ereignisse}
\label{sec:mengenoperationen}

\begin{komplexaufgaben}{sec:mengenoperationen}

	\komplexaufgabe{
		\textbf{Mengenoperationen am Würfelexperiment}

		Beim Werfen eines fairen Würfels ist die Ergebnismenge $\Omega = \{1, 2, 3, 4, 5, 6\}$.
		Betrachten Sie folgende Ereignisse (jeweils als Menge von Ausgängen):
		\[
			A = \{2, 4, 6\},\quad B = \{1, 2, 3\},\quad C = \{3, 6\}.
		\]
		($A$: gerade Augenzahl; $B$: Augenzahl kleiner als 4; $C$: durch 3 teilbar)
	}{
		\item Bestimmen Sie $A \cap B$ und beschreiben Sie das Ergebnis in Worten.

		\item Bestimmen Sie $A \cup C$ und beschreiben Sie das Ergebnis in Worten.

		\item Bestimmen Sie $B^c$ und beschreiben Sie das Ergebnis in Worten.

		\item Bestimmen Sie die Differenz $A \setminus B = \{x \in \Omega \mid x \in A \text{ und } x \notin B\}$.

		\item Überprüfen Sie das erste De Morgansche Gesetz: Berechnen Sie $(A \cup B)^c$ und $A^c \cap B^c$ und vergleichen Sie die Ergebnisse.

		\item Sind die Ereignisse $A$ und $C$ unvereinbar (d.h.\ gilt $A \cap C = \emptyset$)? Begründen Sie.
	}

	\komplexaufgabe{
		\textbf{Ereignisse als Mengen beim dreimaligen Münzwurf}

		Ein fairer Münze wird dreimal geworfen. Mögliche Ausgänge notieren wir als Folge von
		K (Kopf) und Z (Zahl):
		\[
			\Omega = \{\mathrm{KKK},\,\mathrm{KKZ},\,\mathrm{KZK},\,\mathrm{KZZ},\,\mathrm{ZKK},\,\mathrm{ZKZ},\,\mathrm{ZZK},\,\mathrm{ZZZ}\}.
		\]
	}{
		\item Listen Sie die Elemente der folgenden Ereignisse auf:
		\begin{itemize}
			\item $A$ = „mindestens zweimal Kopf"
			\item $B$ = „das erste Ergebnis ist Kopf"
			\item $C$ = „alle drei Ergebnisse sind gleich"
		\end{itemize}

		\item Bestimmen Sie $A \cap B$ und beschreiben Sie dieses Ereignis in Worten.

		\item Bestimmen Sie $A \cup C$ und $B^c$.

		\item Überprüfen Sie das De Morgansche Gesetz $(A \cap B)^c = A^c \cup B^c$, indem Sie beide Seiten durch Aufzählen der Elemente bestimmen und vergleichen.

		\item Bestimmen Sie $P(A)$, $P(B)$ und $P(A \cap B)$ im Laplace-Modell. Gilt $P(A \cup B) = P(A) + P(B) - P(A \cap B)$?
	}

	\komplexaufgabe{
		\textbf{De Morgansche Gesetze und Distributivgesetz nachweisen}

		Sei $\Omega = \{1, 2, 3, 4, 5, 6, 7, 8, 9, 10\}$ sowie
		\[
			A = \{2, 4, 6, 8, 10\},\quad B = \{1, 2, 3, 4, 5\},\quad D = \{4, 5, 6, 7, 8\}.
		\]
	}{
		\item Bestimmen Sie $A^c$, $B^c$, $A \cup B$ und $A \cap B$ durch Aufzählen der Elemente.

		\item Weisen Sie das erste De Morgansche Gesetz nach:
		\[
			(A \cup B)^c = A^c \cap B^c.
		\]

		\item Weisen Sie das zweite De Morgansche Gesetz nach:
		\[
			(A \cap B)^c = A^c \cup B^c.
		\]

		\item Überprüfen Sie ein Distributivgesetz:
		\[
			A \cap (B \cup D) = (A \cap B) \cup (A \cap D).
		\]
	}

	\komplexaufgabe{
		\textbf{Anwendung: Ereignisse im Skatspiel}

		Ein Skatspiel besteht aus 32 Karten: vier Farben (Kreuz, Pik, Herz, Karo), jeweils mit den Werten 7, 8, 9, 10, Bube, Dame, König, As. Eine Karte wird zufällig gezogen.

		Betrachten Sie die Ereignisse:
		\begin{itemize}
			\item $H$ = „die Karte ist eine Herzkarte" \quad ($|H| = 8$)
			\item $B$ = „die Karte ist eine Bildkarte" (Bube, Dame, König oder As) \quad ($|B| = 16$)
			\item $J$ = „die Karte ist ein Bube" \quad ($|J| = 4$)
		\end{itemize}
	}{
		\item Bestimmen Sie $H \cap B$. Wie viele Karten enthält diese Menge?

		\item Bestimmen Sie $|H \cup B|$ mit Hilfe der Formel $|H \cup B| = |H| + |B| - |H \cap B|$.

		\item Bestimmen Sie $H^c$ und beschreiben Sie dieses Gegenereignis in Worten.

		\item Gilt $J \subseteq B$? Begründen Sie mit der Definition der Teilmenge.

		\item Berechnen Sie im Laplace-Modell $P(H)$, $P(B)$, $P(H \cap B)$ und $P(H \cup B)$.
			Überprüfen Sie: $P(H \cup B) = P(H) + P(B) - P(H \cap B)$.
	}

\end{komplexaufgaben}


\section{Stochastik als mathematische Theorie}
\label{sec:stochTheorie}

\begin{komplexaufgaben}{sec:stochTheorie}

	\komplexaufgabe{
		\textbf{Ein diskret-unendlicher Grundraum}

		In einer Notrufzentrale wird jeder eingehende Anruf als „Notfall" oder „kein Notfall"
		eingestuft. Beobachtet wird der Anruf, bei dem erstmals ein Notfall eintritt.
		Der Grundraum ist $\Omega = \{1, 2, 3, \ldots\}$.
	}{
		\item Beschreiben Sie in Worten, was das Ergebnis $5$ bedeutet.
		\item Formulieren Sie die folgenden Ereignisse als Teilmengen von $\Omega$:
		\begin{itemize}
			\item $A$: Der erste Notfall tritt spätestens beim 4.\ Anruf ein.
			\item $B$: Der erste Notfall tritt bei einem geraden Anruf ein.
			\item $C$: Der erste Notfall tritt frühestens beim 3.\ Anruf ein.
		\end{itemize}
		\item Bestimmen Sie $A \cap C$, $A \cup C$ und $B^c$.
		\item Erklären Sie, warum man hier nicht mit „günstige Fälle durch alle Fälle" arbeiten kann.
	}

	\komplexaufgabe{
		\textbf{Wahrscheinlichkeiten in einem diskret-unendlichen Modell}

		Im Modell aus der vorigen Aufgabe sei die Wahrscheinlichkeit dafür, dass ein einzelner
		Anruf ein Notfall ist, gleich $p = 0{,}2$. Für $n = 1, 2, 3, \ldots$ gelte:
		\[
			P(\{n\}) = 0{,}8^{n-1} \cdot 0{,}2.
		\]
	}{
		\item Berechnen Sie $P(\{1\})$, $P(\{2\})$ und $P(\{3\})$.
		\item Bestimmen Sie $P(A)$ für $A = \{1, 2, 3, 4\}$.
		\item Deuten Sie das Ereignis $A$ in Worten.
		\item Vergleichen Sie dieses Modell mit einem Laplace-Modell. Was ist der entscheidende Unterschied?
	}

	\komplexaufgabe{
		\textbf{Einfacher stetiger Grundraum}

		Ein Messfehler werde durch eine Zahl im Intervall $[-0{,}5;\, 0{,}5]$ beschrieben.
		Als Grundraum wird $\Omega = [-0{,}5;\, 0{,}5]$ verwendet. Gegeben seien die Ereignisse:
		\[
			A = [-0{,}1;\, 0{,}1], \quad B = [0;\, 0{,}3].
		\]
	}{
		\item Bestimmen Sie $A \cap B$, $A \cup B$ und $A^c$.
		\item Beschreiben Sie die drei Ereignisse in Worten.
		\item Erklären Sie, warum man hier die Ereignisse nicht durch Abzählen behandeln kann.
		\item Nennen Sie zwei weitere sinnvolle Ereignisse in diesem Modell.
	}

	\komplexaufgabe{
		\textbf{Gleichverteilung auf einem Intervall als Modellidee}

		Im Modell aus der vorigen Aufgabe werde vereinfachend angenommen,
		dass gleich lange Intervalle gleich wahrscheinlich sind.
	}{
		\item Berechnen Sie die Wahrscheinlichkeit von $A = [-0{,}1;\, 0{,}1]$.
		\item Berechnen Sie die Wahrscheinlichkeit von $B = [0;\, 0{,}3]$.
		\item Berechnen Sie die Wahrscheinlichkeit von $A \cap B$.
		\item Überprüfen Sie für diese Ereignisse die Formel
			$P(A \cup B) = P(A) + P(B) - P(A \cap B)$.
	}

	\komplexaufgabe{
		\textbf{Laplace als Spezialfall erkennen}

		Betrachten Sie die drei folgenden Situationen:
		\begin{enumerate}[label=(\roman*)]
			\item einmaliges Werfen eines fairen Würfels,
			\item Warten auf den ersten Notfall (Modell aus den Aufgaben 1 und 2),
			\item Messfehler in einem Intervall (Modell aus den Aufgaben 3 und 4).
		\end{enumerate}
	}{
		\item Geben Sie jeweils einen sinnvollen Grundraum an.
		\item Entscheiden Sie, ob es sich um einen endlichen, diskret-unendlichen oder stetigen
			Grundraum handelt.
		\item Entscheiden Sie, ob das Laplace-Modell jeweils passt, und begründen Sie.
		\item Erklären Sie in einem kurzen Text, warum die Stochastik eine allgemeine Sprache
			für alle drei Situationen braucht.
	}

	\komplexaufgabe{
		\textbf{Folgerungen aus Grundregeln}

		Seien $A$ und $B$ zwei Ereignisse in einem Zufallsexperiment. Es gelte:
		\[
			P(A) = 0{,}6, \quad P(B) = 0{,}4, \quad P(A \cap B) = 0{,}15.
		\]
	}{
		\item Berechnen Sie $P(A \cup B)$.
		\item Berechnen Sie $P(A^c)$.
		\item Entscheiden Sie, ob $A$ und $B$ unvereinbar sind.
		\item Erklären Sie, welche dieser Rechnungen nicht vom konkreten Zufallsexperiment abhängen.
	}

	\komplexaufgabe{
		\textbf{Gegenereignis als Denkwerkzeug}

		Ein Ereignis $A$ bedeute: „Beim dreimaligen Werfen einer fairen Münze fällt
		mindestens einmal Kopf."
	}{
		\item Beschreiben Sie das Gegenereignis $A^c$ in Worten.
		\item Bestimmen Sie $P(A^c)$.
		\item Berechnen Sie daraus $P(A)$.
		\item Erklären Sie, warum das Arbeiten mit dem Gegenereignis oft einfacher ist.
	}

\end{komplexaufgaben}


\setcounter{section}{1}
\setcounter{subsection}{0}
\begin{loesungssektion}{sec:kombinatorik}
	\setcounter{exo}{0}

	\begin{komplexloesungen}{sec:kombinatorik}

	\komplexloesung{}{
		\item $9^6 = 531\,441$
		\item $\binom{20}{10} = 184\,756$
		\item $\frac{8!}{(8-3)!} = 8 \cdot 7 \cdot 6 = 336$
		\item $\frac{8!}{(8-5)!} = 8 \cdot 7 \cdot 6 \cdot 5 \cdot 4 = 6\,720$
		\item $5! = 120$
		\item $3^{11} = 177\,147$
		\item $\binom{8}{5} = 56$
		\item $\binom{34}{3} = 5\,984$
		\item $\binom{19}{11} = 75\,582$
		\item $5^3 = 125$
	}

	\komplexloesung{$P(\text{Fuzzi kommt durch}) = \dfrac{6}{7} \cdot \dfrac{4}{5} \cdot \dfrac{2}{3}
		= \dfrac{48}{105} \approx 45{,}71\,\%$. Da $P < 50\,\%$, sollte man \emph{nicht} wetten.}{}

	\komplexloesung{$2^{8} = 256$, \quad $2^{16} = 65\,536$, \quad $2^{32} = 4\,294\,967\,296$}{}

	\komplexloesung{$\binom{9}{7} = 36$ Sieben-er-Auswahlen $\;\times\; 7! = 5\,040$ Anordnungen
		$= 181\,440$}{}

	\komplexloesung{Person 1 stößt mit den Gästen 2,\,\ldots,\,7 an; Person 2 mit 3,\,\ldots,\,7 usw.\\
		$6+5+4+3+2+1 = 21$ (bzw.\ $\binom{7}{2} = 21$)}{}

	\end{komplexloesungen}

\end{loesungssektion}


\setcounter{section}{2}
\setcounter{subsection}{0}
\begin{loesungssektion}{sec:klassWahrsch}
	\setcounter{exo}{0}

	\begin{komplexloesungen}{sec:klassWahrsch}

	\komplexloesung{}{
		\item $P(\text{genau 3 rote}) = \dfrac{\binom{5}{3}\binom{15}{5}}{\binom{20}{8}}
				\approx 23{,}84\,\%$;
			\quad $P(\text{mind.\,4 rote}) = P(\text{4r}) + P(\text{5r})
				\approx 5{,}42\,\% + 0{,}36\,\% = 5{,}78\,\%$
		\item $P(\text{genau 2 defekt}) = \dfrac{\binom{10}{2}\binom{90}{3}}{\binom{100}{5}}
				\approx 7{,}02\,\%$;
			\quad $P(\text{mind.\,3 defekt}) \approx 0{,}64\,\%+0{,}03\,\%+0{,}0003\,\% \approx 0{,}66\,\%$
	}

	\komplexloesung{$P(\text{3 Richtige}) \approx 0{,}0264$, \;
		$P(\text{4 Richtige}) \approx 0{,}000549$.\\
		Mittlerer Gewinn des Veranstalters: $G \approx 1 - 0{,}0264 \cdot 10 - 0{,}000549 \cdot 1000
		\approx 0{,}19\,\text{€}$ pro Tipp}{}

	\komplexloesung{$P(\text{alle verschiedenfarbig}) =
		\dfrac{\binom{5}{1}\binom{3}{1}\binom{6}{1}}{\binom{14}{3}} = \dfrac{90}{364} \approx 24{,}7\,\%$;
		\quad $P(\text{alle rot}) = \dfrac{\binom{5}{3}}{\binom{14}{3}} = \dfrac{10}{364} \approx 2{,}7\,\%$;
		\quad $P(\text{alle gleichfarbig}) = \dfrac{10+1+20}{364} = \dfrac{31}{364} \approx 8{,}5\,\%$}{}

	\komplexloesung{}{
		\item $P(\text{genau 1 defekt}) = \dfrac{\binom{10}{1}\binom{70}{4}}{\binom{80}{5}}
				\approx 0{,}3814$
		\item $P(\text{genau 3 defekt}) = \dfrac{\binom{10}{3}\binom{70}{2}}{\binom{80}{5}}
				\approx 0{,}0121$
		\item $P(\text{höchstens 4 defekt}) = 1 - \dfrac{\binom{10}{5}\binom{70}{0}}{\binom{80}{5}}
				\approx 0{,}9999895$
		\item $P(\text{mindestens 1 defekt}) = 1 - \dfrac{\binom{10}{0}\binom{70}{5}}{\binom{80}{5}}
				\approx 0{,}4965$
	}

	\komplexloesung{}{
		\item $P(\text{1 HG, 4 Nieten}) = \dfrac{\binom{2}{1}\binom{70}{4}}{\binom{100}{5}}
				\approx 0{,}0244$;\quad
			$P(\text{kein Gewinn}) = \dfrac{\binom{70}{5}}{\binom{100}{5}} \approx 0{,}1608$
		\item $P(\text{2 EG, 3 TP, 5 N}) = \dfrac{\binom{8}{2}\binom{20}{3}\binom{70}{5}}{\binom{100}{10}}
				\approx 0{,}0223$;\quad
			$P(\text{1 HG, 2 EG, 7 N}) = \dfrac{\binom{2}{1}\binom{8}{2}\binom{70}{7}}{\binom{100}{10}}
				\approx 0{,}0039$
	}

	\komplexloesung{\textbf{Das Problem des Chevalier de Méré}}{
		\item \emph{(Lernendendiskussion — keine eindeutige Musterlösung)}

		\item Wette~A: $\Omega = \{1,\ldots,6\}^4$, Ereignis: „mindestens eine 6". \\
		      Wette~B: zwei Würfel, 24 Würfe, Ereignis: „mindestens eine Doppel-6". \\
		      Annahme: alle Würfe unabhängig, faire Würfel.

		\item $P(\text{keine 6 in 4 Würfen}) = \left(\dfrac{5}{6}\right)^4 \approx 0{,}482$; \quad
		      $P(\text{mind. eine 6}) = 1 - \left(\dfrac{5}{6}\right)^4 \approx 0{,}518$

		\item $P(\text{Doppel-6 bei einem Wurf}) = \dfrac{1}{36}$; \quad
		      $P(\text{mind. eine Doppel-6 in 24 Würfen}) = 1 - \left(\dfrac{35}{36}\right)^{24} \approx 0{,}491$

		\item Wette~A ist günstiger (${\approx}\,51{,}8\,\%$ gegenüber ${\approx}\,49{,}1\,\%$).

		\item De Méré hat Wahrscheinlichkeiten addiert statt die Gegenwahrscheinlichkeit zu benutzen; korrekt ist $P = 1-(1-p)^n$.

		\item Der Pascal-Fermat-Briefwechsel (1654) begründete die systematische Behandlung von Zufallsexperimenten und gilt als Geburtsstunde der Wahrscheinlichkeitsrechnung.
	}

	% --- Lösungen Bernoulli-Grundaufgaben ---

	\komplexloesung{Würfel 4-mal, $p=\tfrac{1}{6}$}{
		\item Baumdiagramm: Erste Stufe — zwei Äste: „6" (Wahrsch.\ $\tfrac{1}{6}$) und „nicht 6" (Wahrsch.\ $\tfrac{5}{6}$). Zweite Stufe analog.
		\item $P(\text{genau 1}) = \binom{4}{1}\cdot\dfrac{1}{6}\cdot\left(\dfrac{5}{6}\right)^3
			= 4 \cdot \dfrac{125}{1296} = \dfrac{500}{1296} \approx 38{,}6\,\%$
		\item $P(\text{mindestens 1}) = 1 - \left(\dfrac{5}{6}\right)^4
			= 1 - \dfrac{625}{1296} = \dfrac{671}{1296} \approx 51{,}8\,\%$
		\item $P(\text{keine 6}) = \left(\dfrac{5}{6}\right)^4 = \dfrac{625}{1296} \approx 48{,}2\,\%$
	}

	\komplexloesung{Münze 5-mal, $p=\tfrac{1}{2}$}{
		\item $P(\text{genau 3}) = \binom{5}{3}\cdot\left(\dfrac{1}{2}\right)^5 = \dfrac{10}{32} = \dfrac{5}{16} \approx 31{,}3\,\%$
		\item $P(\text{höchstens 2}) = 1 - P(3) - P(4) - P(5)
			= 1 - \dfrac{10}{32} - \dfrac{5}{32} - \dfrac{1}{32} = \dfrac{16}{32} = 50{,}0\,\%$
		\item $P(\text{mindestens 4}) = \dfrac{5}{32} + \dfrac{1}{32} = \dfrac{6}{32} = \dfrac{3}{16} \approx 18{,}75\,\%$
	}

	\komplexloesung{Freiwurf $p=0{,}7$, 3 Würfe}{
		\item Baumdiagramm: je zwei Äste T ($p=0{,}7$) und F ($p=0{,}3$); 8 Pfade.
		\item $P(\text{genau 2}) = \binom{3}{2}\cdot 0{,}7^2 \cdot 0{,}3
			= 3 \cdot 0{,}49 \cdot 0{,}3 = 0{,}441$
		\item $P(\text{mindestens 1}) = 1 - 0{,}3^3 = 1 - 0{,}027 = 0{,}973$
	}

	\komplexloesung{Glühbirnen $p=0{,}05$, $n=6$}{
		\item $P(\text{keine defekte}) = 0{,}95^6 \approx 0{,}7351$
		\item $P(\text{genau 1}) = \binom{6}{1}\cdot 0{,}05\cdot 0{,}95^5
			\approx 6 \cdot 0{,}05 \cdot 0{,}7738 \approx 0{,}2321$
		\item $P(\text{mindestens 1}) = 1 - 0{,}95^6 \approx 0{,}2649$
	}

	\komplexloesung{Multiple-Choice $p=\tfrac{1}{3}$, $n=4$}{
		\item $P(\text{alle 4}) = \left(\dfrac{1}{3}\right)^4 = \dfrac{1}{81} \approx 1{,}2\,\%$
		\item $P(\text{genau 2}) = \binom{4}{2}\cdot\left(\dfrac{1}{3}\right)^2\cdot\left(\dfrac{2}{3}\right)^2
			= 6\cdot\dfrac{4}{81} = \dfrac{24}{81} \approx 29{,}6\,\%$
		\item $P(\text{mindestens 3}) = \binom{4}{3}\cdot\left(\dfrac{1}{3}\right)^3\cdot\dfrac{2}{3}
			+ \left(\dfrac{1}{3}\right)^4
			= \dfrac{8}{81} + \dfrac{1}{81} = \dfrac{9}{81} = \dfrac{1}{9} \approx 11{,}1\,\%$
	}

	\komplexloesung{Schützin $p=0{,}8$, $n=5$}{
		\item $P(\text{genau 4}) = \binom{5}{4}\cdot 0{,}8^4 \cdot 0{,}2
			= 5 \cdot 0{,}4096 \cdot 0{,}2 = 0{,}4096$
		\item $P(\text{mindestens 4}) = 0{,}4096 + 0{,}8^5 = 0{,}4096 + 0{,}3277 \approx 0{,}7373$
		\item $P(\text{alle 5}) = 0{,}8^5 \approx 0{,}3277$. \\
			Vergleich: Genau 4 Treffer (${\approx}\,41\,\%$) ist am wahrscheinlichsten,
			dann alle 5 (${\approx}\,33\,\%$), dann mindestens 4 (${\approx}\,74\,\%$ als Summe).
	}

	% --- Lösungen Parametervariation: unbekanntes n ---

	\komplexloesung{Würfel $n$-mal, $p=\tfrac{1}{6}$}{
		\item $P(\text{mindestens eine 6}) = 1 - \left(\dfrac{5}{6}\right)^n$
		\item Gesucht: $\left(\dfrac{5}{6}\right)^n \le 0{,}10$. \\
			Logarithmieren: $n\cdot\ln\dfrac{5}{6} \le \ln 0{,}10$; da $\ln\tfrac{5}{6} < 0$
			kehrt sich das Zeichen um: $n \ge \dfrac{\ln 0{,}10}{\ln\tfrac{5}{6}}
			\approx \dfrac{-2{,}303}{-0{,}182} \approx 12{,}6$. \\
			Mindestens $n = 13$ Würfe. \;
			Probe: $\left(\tfrac{5}{6}\right)^{13} \approx 0{,}095 < 0{,}10$. $\checkmark$
	}

	\komplexloesung{Münze $n$-mal, mindestens 3 Kopf}{
		\item $n=4$: $P = \binom{4}{3}\cdot\tfrac{1}{16} + \tfrac{1}{16} = \tfrac{5}{16} \approx 31{,}25\,\%$. \quad
			$n=5$: $P = \tfrac{10}{32}+\tfrac{5}{32}+\tfrac{1}{32} = \tfrac{16}{32} = 50{,}0\,\%$.
		\item Ab $n = 5$; für $n=5$ gilt bereits Gleichheit.
	}

	\komplexloesung{Schrauben $p=0{,}1$}{
		\item $P(\text{mindestens 1 defekt}) = 1 - 0{,}9^n$
		\item Gesucht: $0{,}9^n \le 0{,}10$. \\
			$n \ge \dfrac{\ln 0{,}10}{\ln 0{,}9} \approx \dfrac{-2{,}303}{-0{,}105} \approx 21{,}9$. \\
			Mindestens $n = 22$ Schrauben. \;
			Probe: $0{,}9^{22} \approx 0{,}098 < 0{,}10$. $\checkmark$
	}

	% --- Lösungen Parametervariation: unbekanntes p ---

	\komplexloesung{$n=4$, $P(\text{genau 2}) = \tfrac{216}{625}$}{
		\item $\binom{4}{2}\cdot p^2(1-p)^2 = \dfrac{216}{625}$; \quad
			$6\cdot[p(1-p)]^2 = \dfrac{216}{625}$; \quad
			$[p(1-p)]^2 = \dfrac{36}{625}$; \quad
			$p(1-p) = \dfrac{6}{25}$; \quad
			$p^2 - p + \dfrac{6}{25} = 0$; \quad
			$p = \dfrac{1 \pm \sqrt{1 - \tfrac{24}{25}}}{2} = \dfrac{1 \pm \tfrac{1}{5}}{2}$. \\
			Lösungen: $p = \dfrac{3}{5} = 0{,}6$ \; oder \; $p = \dfrac{2}{5} = 0{,}4$.
		\item Symmetrie: Vertauscht man Treffer und Niete (ersetzt $p$ durch $1-p$), bleibt $P(\text{genau 2})$ unverändert.
	}

	\komplexloesung{$n=2$, $P(\text{mindestens 1}) = \tfrac{5}{9}$: \quad
		$1 - (1-p)^2 = \dfrac{5}{9}$; \quad $(1-p)^2 = \dfrac{4}{9}$; \quad
		$1-p = \dfrac{2}{3}$ (da $p \in [0,1]$); \quad $p = \dfrac{1}{3}$.}{}

	\komplexloesung{$n=3$, $P(\text{kein Treffer}) = \tfrac{8}{27}$: \quad
		$(1-p)^3 = \dfrac{8}{27}$; \quad $1-p = \dfrac{2}{3}$; \quad $p = \dfrac{1}{3}$.}{}

	% --- Lösungen Parametervariation: unbekanntes k ---

	\komplexloesung{Würfel 6-mal, $p=\tfrac{1}{6}$, häufigste Anzahl Sechsen: \;
		$k{=}0$: ${\approx}\,33{,}5\,\%$; \;
		$k{=}1$: ${\approx}\,40{,}2\,\%$; \;
		$k{=}2$: ${\approx}\,20{,}1\,\%$; \;
		$k{=}3$: ${\approx}\,5{,}4\,\%$. \;
		Für $k \ge 4$ noch kleiner. Häufigster Wert: $\boldsymbol{k = 1}$.}{}

	\komplexloesung{Münze 10-mal, $P(k) = \binom{10}{k}/1024$: \;
		$k{=}4$: $210/1024 \approx 20{,}5\,\%$ $\checkmark$; \;
		$k{=}5$: $252/1024 \approx 24{,}6\,\%$ $\checkmark$; \;
		$k{=}6$: $210/1024 \approx 20{,}5\,\%$ $\checkmark$. \;
		$k{=}3$ und $k{=}7$: ${\approx}\,11{,}7\,\% < 20\,\%$. \;
		Ergebnis: $k \in \{4, 5, 6\}$.}{}

	\komplexloesung{$n=5$, $p=0{,}3$, kleinstes $k$ mit $P(X \ge k) \le 0{,}10$: \;
		$P(X \ge 3) \approx 0{,}132 + 0{,}028 + 0{,}002 = 0{,}163 > 0{,}10$; \;
		$P(X \ge 4) \approx 0{,}028 + 0{,}002 = 0{,}031 \le 0{,}10$. $\checkmark$ \;
		Kleinstes $k = \mathbf{4}$.}{}

	\end{komplexloesungen}

\end{loesungssektion}


\setcounter{section}{3}
\setcounter{subsection}{0}
\begin{loesungssektion}{sec:mengenoperationen}
	\setcounter{exo}{0}

	\begin{komplexloesungen}{sec:mengenoperationen}

	\komplexloesung{
		\textbf{Mengenoperationen am Würfelexperiment}
	}{
		\item $A \cap B = \{2\}$ — „die Augenzahl ist gerade und kleiner als 4".

		\item $A \cup C = \{2, 3, 4, 6\}$ — „die Augenzahl ist gerade oder durch 3 teilbar".

		\item $B^c = \{4, 5, 6\}$ — „die Augenzahl ist mindestens 4".

		\item $A \setminus B = \{4, 6\}$ — gerade Zahlen, die nicht kleiner als 4 sind.

		\item $(A \cup B)^c$: $A \cup B = \{1,2,3,4,6\}$, also $(A \cup B)^c = \{5\}$. \\
		$A^c \cap B^c = \{1,3,5\} \cap \{4,5,6\} = \{5\}$. Beide Seiten stimmen überein. $\checkmark$

		\item $A \cap C = \{6\} \neq \emptyset$, also sind $A$ und $C$ \emph{nicht} unvereinbar.
	}

	\komplexloesung{
		\textbf{Dreimaliger Münzwurf}
	}{
		\item $A = \{\mathrm{KKK}, \mathrm{KKZ}, \mathrm{KZK}, \mathrm{ZKK}\}$; \quad
		$B = \{\mathrm{KKK}, \mathrm{KKZ}, \mathrm{KZK}, \mathrm{KZZ}\}$; \quad
		$C = \{\mathrm{KKK}, \mathrm{ZZZ}\}$.

		\item $A \cap B = \{\mathrm{KKK}, \mathrm{KKZ}, \mathrm{KZK}\}$ — „mindestens zweimal Kopf und das erste Ergebnis ist Kopf".

		\item $A \cup C = \{\mathrm{KKK}, \mathrm{KKZ}, \mathrm{KZK}, \mathrm{ZKK}, \mathrm{ZZZ}\}$; \quad
		$B^c = \{\mathrm{ZKK}, \mathrm{ZKZ}, \mathrm{ZZK}, \mathrm{ZZZ}\}$.

		\item $(A \cap B)^c = \{\mathrm{KZZ}, \mathrm{ZKK}, \mathrm{ZKZ}, \mathrm{ZZK}, \mathrm{ZZZ}\}$. \\
		$A^c = \{\mathrm{KZZ}, \mathrm{ZKZ}, \mathrm{ZZK}, \mathrm{ZZZ}\}$; $B^c = \{\mathrm{ZKK}, \mathrm{ZKZ}, \mathrm{ZZK}, \mathrm{ZZZ}\}$; \\
		$A^c \cup B^c = \{\mathrm{KZZ}, \mathrm{ZKK}, \mathrm{ZKZ}, \mathrm{ZZK}, \mathrm{ZZZ}\}$. Beide Seiten gleich. $\checkmark$

		\item $P(A) = \tfrac{4}{8} = \tfrac{1}{2}$; $P(B) = \tfrac{4}{8} = \tfrac{1}{2}$; $P(A \cap B) = \tfrac{3}{8}$. \\
		$P(A \cup B) = \tfrac{5}{8}$; \; $P(A) + P(B) - P(A \cap B) = \tfrac{1}{2} + \tfrac{1}{2} - \tfrac{3}{8} = \tfrac{5}{8}$. $\checkmark$
	}

	\komplexloesung{
		\textbf{De Morgansche Gesetze und Distributivgesetz}
	}{
		\item $A^c = \{1,3,5,7,9\}$; $B^c = \{6,7,8,9,10\}$; $A \cup B = \{1,2,3,4,5,6,8,10\}$; $A \cap B = \{2,4\}$.

		\item $(A \cup B)^c = \{7,9\}$; \; $A^c \cap B^c = \{1,3,5,7,9\} \cap \{6,7,8,9,10\} = \{7,9\}$. Gleich. $\checkmark$

		\item $(A \cap B)^c = \{1,3,5,6,7,8,9,10\}$; \;
		$A^c \cup B^c = \{1,3,5,7,9\} \cup \{6,7,8,9,10\} = \{1,3,5,6,7,8,9,10\}$. Gleich. $\checkmark$

		\item $B \cup D = \{1,2,3,4,5,6,7,8\}$; \; $A \cap (B \cup D) = \{2,4,6,8\}$. \\
		$(A \cap B) = \{2,4\}$; $(A \cap D) = \{4,6,8\}$; $(A \cap B) \cup (A \cap D) = \{2,4,6,8\}$. Gleich. $\checkmark$
	}

	\komplexloesung{
		\textbf{Skatspiel}
	}{
		\item $H \cap B$: die vier Herz-Bildkarten (Herz-Bube, -Dame, -König, -As). $|H \cap B| = 4$.

		\item $|H \cup B| = 8 + 16 - 4 = 20$.

		\item $H^c$: alle Karten, die \emph{keine} Herzkarten sind (Kreuz, Pik, Karo). $|H^c| = 24$.

		\item Ja, $J \subseteq B$: jeder Bube ist eine Bildkarte, also ist jedes Element von $J$ auch in $B$.

		\item $P(H) = \tfrac{8}{32} = \tfrac{1}{4}$; $P(B) = \tfrac{16}{32} = \tfrac{1}{2}$; $P(H \cap B) = \tfrac{4}{32} = \tfrac{1}{8}$; \\
		$P(H \cup B) = \tfrac{20}{32} = \tfrac{5}{8}$; \;
		$P(H) + P(B) - P(H \cap B) = \tfrac{1}{4} + \tfrac{1}{2} - \tfrac{1}{8} = \tfrac{5}{8}$. $\checkmark$
	}

	\end{komplexloesungen}

\end{loesungssektion}


\setcounter{section}{4}
\setcounter{subsection}{0}
\begin{loesungssektion}{sec:stochTheorie}
	\setcounter{exo}{0}

	\begin{komplexloesungen}{sec:stochTheorie}

	\komplexloesung{
		\textbf{Diskret-unendlicher Grundraum}
	}{
		\item Ergebnis $5$: Der erste Notfall tritt beim 5.\ Anruf ein; die ersten vier Anrufe
			waren kein Notfall.

		\item $A = \{1, 2, 3, 4\}$; \quad
			$B = \{2, 4, 6, 8, \ldots\}$ (alle geraden natürlichen Zahlen); \quad
			$C = \{3, 4, 5, 6, \ldots\}$.

		\item $A \cap C = \{3, 4\}$; \quad
			$A \cup C = \{1, 2, 3, 4, 5, \ldots\} = \Omega$; \quad
			$B^c = \{1, 3, 5, 7, \ldots\}$ (ungerade Zahlen).

		\item Das Laplace-Modell setzt voraus, dass alle Ergebnisse gleichwahrscheinlich sind und
			die Ergebnismenge endlich ist. Hier ist $\Omega$ unendlich; außerdem sind die
			Ergebnisse nicht gleichwahrscheinlich.
	}

	\komplexloesung{
		\textbf{Wahrscheinlichkeiten im diskret-unendlichen Modell}
	}{
		\item $P(\{1\}) = 0{,}8^0 \cdot 0{,}2 = 0{,}2$; \quad
			$P(\{2\}) = 0{,}8^1 \cdot 0{,}2 = 0{,}16$; \quad
			$P(\{3\}) = 0{,}8^2 \cdot 0{,}2 = 0{,}128$.

		\item $P(A) = P(\{1\}) + P(\{2\}) + P(\{3\}) + P(\{4\})$\\
			$= 0{,}2 + 0{,}16 + 0{,}128 + 0{,}8^3 \cdot 0{,}2
			= 0{,}2 + 0{,}16 + 0{,}128 + 0{,}1024 = 0{,}5904$.

		\item $A = \{1, 2, 3, 4\}$: Der erste Notfall tritt spätestens beim 4.\ Anruf ein.

		\item Im Laplace-Modell wären alle Ergebnisse gleichwahrscheinlich; hier nimmt
			$P(\{n\})$ mit wachsendem $n$ streng ab. Außerdem setzt das Laplace-Modell eine
			endliche Ergebnismenge voraus — hier ist $\Omega$ unendlich.
	}

	\komplexloesung{
		\textbf{Stetiger Grundraum}
	}{
		\item $A \cap B = [0;\, 0{,}1]$; \quad
			$A \cup B = [-0{,}1;\, 0{,}3]$; \quad
			$A^c = [-0{,}5;\, -0{,}1) \cup (0{,}1;\, 0{,}5]$.

		\item $A$: Der Messfehler ist höchstens $0{,}1$ dem Betrag nach (kleiner Fehler). \\
			$B$: Der Messfehler ist nicht negativ und höchstens $0{,}3$. \\
			$A^c$: Der Betrag des Messfehlers übersteigt $0{,}1$.

		\item $\Omega = [-0{,}5;\, 0{,}5]$ ist überabzählbar unendlich; es gibt keine vollständige
			Liste aller Ergebnisse. Abzählen setzt eine endliche (oder abzählbare)
			Ergebnismenge voraus.

		\item Zum Beispiel: $D = [0{,}2;\, 0{,}4]$ (großer positiver Fehler) oder
			$E = [-0{,}5;\, 0]$ (negativer Messfehler).
	}

	\komplexloesung{
		\textbf{Gleichverteilung auf einem Intervall}
	}{
		\item $P(A) = \dfrac{\text{Länge}(A)}{\text{Länge}(\Omega)} = \dfrac{0{,}2}{1} = 0{,}2$.

		\item $P(B) = \dfrac{0{,}3}{1} = 0{,}3$.

		\item $P(A \cap B) = P([0;\, 0{,}1]) = \dfrac{0{,}1}{1} = 0{,}1$.

		\item $P(A) + P(B) - P(A \cap B) = 0{,}2 + 0{,}3 - 0{,}1 = 0{,}4$. \\
			Direktprobe: $A \cup B = [-0{,}1;\, 0{,}3]$ hat Länge $0{,}4$,
			also $P(A \cup B) = 0{,}4$. $\checkmark$
	}

	\komplexloesung{
		\textbf{Laplace als Spezialfall}
	}{
		\item (i) $\Omega = \{1, 2, 3, 4, 5, 6\}$; \quad
			(ii) $\Omega = \{1, 2, 3, \ldots\}$; \quad
			(iii) $\Omega = [-0{,}5;\, 0{,}5]$.

		\item (i) endlich; \quad (ii) diskret-unendlich; \quad (iii) stetig (überabzählbar unendlich).

		\item (i) Ja: 6 gleichwahrscheinliche Ergebnisse, Laplace anwendbar. \\
			(ii) Nein: $\Omega$ ist unendlich, die Ergebnisse sind nicht gleichwahrscheinlich. \\
			(iii) Nein: $\Omega$ ist überabzählbar; einzelne Punkte haben Wahrscheinlichkeit $0$.

		\item Die drei Situationen haben grundlegend verschiedene Typen von Grundräumen.
			Eine allgemeine Theorie (nach Kolmogorov) beschreibt Ereignisse einheitlich als
			Teilmengen von $\Omega$ und legt Wahrscheinlichkeiten durch Axiome fest —
			unabhängig davon, ob $\Omega$ endlich, diskret-unendlich oder stetig ist.
			Der Laplace-Fall ist dabei nur ein besonders einfacher Spezialfall.
	}

	\komplexloesung{
		\textbf{Folgerungen aus Grundregeln}
	}{
		\item $P(A \cup B) = P(A) + P(B) - P(A \cap B) = 0{,}6 + 0{,}4 - 0{,}15 = 0{,}85$.

		\item $P(A^c) = 1 - P(A) = 1 - 0{,}6 = 0{,}4$.

		\item $A$ und $B$ sind \emph{nicht} unvereinbar: $P(A \cap B) = 0{,}15 \neq 0$,
			also gilt $A \cap B \neq \emptyset$.

		\item Alle vier Rechnungen folgen ausschließlich aus den allgemeinen
			Wahrscheinlichkeitsregeln (Additionssatz, Komplementregel) und gelten in
			\emph{jedem} Wahrscheinlichkeitsraum — unabhängig vom konkreten Experiment.
	}

	\komplexloesung{
		\textbf{Gegenereignis als Denkwerkzeug}
	}{
		\item $A^c$ = „Beim dreimaligen Werfen einer fairen Münze fällt kein einziges Mal Kopf"
			(also dreimal Zahl, d.\,h.\ $A^c = \{\mathrm{ZZZ}\}$).

		\item $P(A^c) = \left(\dfrac{1}{2}\right)^3 = \dfrac{1}{8}$.

		\item $P(A) = 1 - P(A^c) = 1 - \dfrac{1}{8} = \dfrac{7}{8}$.

		\item Das Gegenereignis „kein einziges Mal Kopf" hat eine viel einfachere Struktur als
			„mindestens einmal Kopf": Es gibt nur ein einziges günstiges Ergebnis (ZZZ), während
			für $A$ sieben Ergebnisse günstig wären. Immer wenn „mindestens einmal" gefragt ist,
			führt der Weg über das Gegenereignis schneller zum Ziel.
	}

	\end{komplexloesungen}

\end{loesungssektion}

\end{document}
